\documentclass[11pt,a4paper]{article}
\usepackage[margin=2.5cm]{geometry}
\usepackage{amsmath,amssymb}

\begin{document}

\section*{Homework 1 Group I}
Member names to be added.

\section*{Question 1}

Let \(s_i=1-2y_i\). For \(\lambda>0\), the objective is
\[
f_\lambda(\theta)
=
\underbrace{\sum_{i=1}^{m}
\phi\!\left(s_i\langle\widetilde{x}_i,\theta\rangle\right)}_{f(\theta)}
+\frac{\lambda}{2}\|\theta\|^2,
\qquad \theta\in E.
\]
We first show that \(f\) is convex. The smoothed hinge loss satisfies
\[
\phi'(z)=
\begin{cases}
0, & z\leq -1,\\
1+z, & -1<z<0,\\
1, & z\geq 0.
\end{cases}
\]
Indeed, the values and one-sided derivatives of the defining pieces agree
at \(z=-1\) and \(z=0\). Thus \(\phi\) is continuously differentiable.
Its derivative is nondecreasing, so for any \(a,b\in\mathbb R\),
\[
\phi(b)-\phi(a)-\phi'(a)(b-a)
=\int_a^b \bigl(\phi'(u)-\phi'(a)\bigr)\,du\geq 0.
\]
(The inequality holds for either ordering of \(a\) and \(b\).)
By Theorem~4.21 of the lecture notes, \(\phi\) is convex.

For any \(\theta,\eta\in E\) and \(t\in[0,1]\), linearity of the inner
product and convexity of \(\phi\) give
\begin{align*}
\phi\!\left(s_i\langle\widetilde{x}_i,(1-t)\theta+t\eta\rangle\right)
&=\phi\!\left((1-t)s_i\langle\widetilde{x}_i,\theta\rangle
+t s_i\langle\widetilde{x}_i,\eta\rangle\right)\\
&\leq (1-t)\phi\!\left(s_i\langle\widetilde{x}_i,\theta\rangle\right)
+t\phi\!\left(s_i\langle\widetilde{x}_i,\eta\rangle\right).
\end{align*}
Summing over \(i\) yields
\(f((1-t)\theta+t\eta)\leq(1-t)f(\theta)+tf(\eta)\);
hence \(f\) is convex.

Now
\[
f_\lambda(\theta)-\frac{\lambda}{2}\|\theta\|^2=f(\theta)
\]
is convex. Since \(\lambda>0\), Definition~4.5 of the lecture notes
therefore implies that \(f_\lambda\) is \(\lambda\)-strongly convex.
It is real-valued on the whole Euclidean space \(E\), so
Theorem~4.28 applies: \(f_\lambda\) has a unique global minimizer
\(\theta^\star\).

Finally, \(f_\lambda\) is convex and differentiable on \(E\):
each loss term is a composition of the continuously differentiable
function \(\phi\) with a linear map, and the quadratic term is
differentiable. By Corollary~4.23 of the lecture notes, a point is
a global minimizer if and only if its gradient vanishes. Consequently,
\[
\{\theta\in E:\nabla f_\lambda(\theta)=0\}=\{\theta^\star\}.
\]

\section*{Question 2}

We now set \(\lambda=0.005\). Write
\(g_i(\theta)=s_i\langle\widetilde{x}_i,\theta\rangle\).
Each \(g_i:E\to\mathbb R\) is linear, hence differentiable, and
\(\phi\) is continuously differentiable as shown in Question 1.
Thus \(f=\sum_{i=1}^m\phi\circ g_i\) is differentiable.
By additivity of the differential and the chain rule, for any \(h\in E\),
\[
Df(\theta)[h]
=\sum_{i=1}^m \phi'(g_i(\theta))\,Dg_i(\theta)[h].
\]
Using the limit expression for the differential of \(g_i\), we obtain
\[
Dg_i(\theta)[h]
=\lim_{t\to0}\frac{g_i(\theta+th)-g_i(\theta)}{t}
=\lim_{t\to0}\frac{t s_i\langle\widetilde{x}_i,h\rangle}{t}
=s_i\langle\widetilde{x}_i,h\rangle.
\]
Substituting and using bilinearity of the inner product gives
\begin{align*}
Df(\theta)[h]
&=\sum_{i=1}^m
s_i\phi'\!\left(s_i\langle\widetilde{x}_i,\theta\rangle\right)
\langle\widetilde{x}_i,h\rangle\\
&=\left\langle
\sum_{i=1}^m s_i\phi'\!\left(s_i\langle\widetilde{x}_i,\theta\rangle\right)
\widetilde{x}_i,\ h
\right\rangle.
\end{align*}

The quadratic term \(r(\theta)=\frac{\lambda}{2}\|\theta\|^2\)
is also differentiable. Indeed,
\[
r(\theta+h)-r(\theta)-\lambda\langle\theta,h\rangle
=\frac{\lambda}{2}\|h\|^2=o(\|h\|),
\]
so \(Dr(\theta)[h]=\lambda\langle\theta,h\rangle\).
Consequently, \(f_\lambda=f+r\) is differentiable and
\[
Df_\lambda(\theta)[h]
=\left\langle
\sum_{i=1}^m s_i\phi'\!\left(s_i\langle\widetilde{x}_i,\theta\rangle\right)
\widetilde{x}_i+\lambda\theta,\ h
\right\rangle.
\]
By Definition~2.18 of the lecture notes,
\(Df_\lambda(\theta)[h]=\langle\nabla f_\lambda(\theta),h\rangle\).
Since this equality holds for every \(h\in E\), we conclude that
\[
\boxed{
\nabla f_\lambda(\theta)
=\sum_{i=1}^m
s_i\phi'\!\left(s_i\langle\widetilde{x}_i,\theta\rangle\right)
\widetilde{x}_i+\lambda\theta.
}
\]


\end{document}
